\nonstopmode 
%\documentclass[11pt,a4paper,twoside]{tesis}
% SI NO PENSAS IMPRIMIRLO EN FORMATO LIBRO PODES USAR
\documentclass[11pt,a4paper]{tesis}
\usepackage{titlesec}
\titleformat{\chapter}[hang]{\normalfont\huge\bfseries}{\thechapter}{1em}{}
\titlespacing*{\chapter}{0pt}{0pt}{20pt}

\makeatletter
\renewcommand\chapter{\if@openright\else\fi
  \thispagestyle{plain}%
  \global\@topnum\z@
  \@afterindentfalse
  \secdef\@chapter\@schapter}
\makeatother

\usepackage{graphicx}
\usepackage[utf8]{inputenc}
\usepackage[spanish]{babel}
\usepackage[left=3cm,right=3cm,bottom=3.5cm,top=3.5cm]{geometry}

\begin{document}

%%%% CARATULA

\def\autor{Lautaro Lasorsa}
\def\tituloTesis{Propuesta de mejora en la medición del bienestar económico}
\def\runtitulo{Propuesta de mejora en la medición del bienestar económico}
\def\director{Rodrigo Castro}
\def\codirector{Walter Sosa Escudero}
\def\lugar{Buenos Aires, Argentina, 2024}
\newcommand{\HRule}{\rule{\linewidth}{0.2mm}}
%
\thispagestyle{empty}

\begin{center}\leavevmode

\vspace{-2cm}

\begin{tabular}{l}
\includegraphics[width=2.6cm]{logofcen.pdf}
\end{tabular}


{\large \sc Universidad de Buenos Aires

Facultad de Ciencias Exactas y Naturales}

\vspace{7.0cm}

%\vspace{3.0cm}
%{
%\Large \color{red}
%\begin{tabular}{|p{2cm}cp{2cm}|}
%\hline
%& Pre-Final Version: \today &\\
%\hline
%\end{tabular}
%}
%\vspace{2.5cm}

\begin{huge}
\textbf{\tituloTesis}
\end{huge}

\vspace{2cm}

{\large Tesis de Licenciatura en Ciencias de Datos}

\vspace{2cm}

{\Large \autor}

\end{center}

\vfill

{\large

{Director: \director}

\vspace{.2cm}

{Codirector: \codirector}

\vspace{.2cm}

\lugar
}

\newpage\thispagestyle{empty}


%%%% ABSTRACTS, AGRADECIMIENTOS Y DEDICATORIA
% \frontmatter

% \tableofcontents

% \mainmatter
% \pagestyle{headings}

%%%% ACA VA EL CONTENIDO DE LA TESIS

\chapter{Medición del Bienestar Económico}

En este capítulo se presenta la problemática a tratar en la tesis, se motiva su estudio y se comparan las formas actuales de medirlo con las alternativas propuestas.

\section{Definición y motivación}

En esta sección definimos qué es el bienestar económico y otros conceptos relacionados y por qué es importante medirlo bien.

\section{Medición actual}

En esta sección comentamos la forma actualmente utilizada, entre otros por quienes elaboran IDH, para medir el bienestar económico promedio de una población, analizando sus ventajas y desventajas.

\section{Propuesta de mejora: Medición ideal}

En esta sección comentaremos cuál es la que proponemos como la forma ideal de medir el bienestar económico, la justificación de esta propuesta y sus limitaciones, sobre todo a nivel logístico.

\section{Propuesta de mejora: Mediciones aproximadas}

En esta sección introducimos el concepto de granularidad de una medición, y estudiamos soluciones de compromiso para crear una medición que se aproxima más a la que proponemos como medición ideal pero siendo logísticamente factibles.
Serán las mediciones que utilizaremos con los datos empiricos.

\chapter{Técnicas a utilizar}

En este capítulo se explicarán brevemente las métricas y técnicas más importantes a utilizar en la presente tesis.

\section{Correlación de Pearson}

Coeficiente clásico de correlación lineal

\section{Correlaciones no lineales}

Coeficientes de correlación de Spearman y $\tau$ de Kendall.

\section{Bootstrap}

Técnica de remuestreo para estimar la distribución de los datos o de una función de los datos.

\chapter{Datos sintéticos}

En este capítulo se utilizan datos simulados para explorar características de los indicadores (existentes y propuestos) en sí mismos.

\section{Generación}
En esta sección se comentará qué datos sintéticos se generaron, qué asunciones se hicieron sobre la distribución de los mismos, y un breve comentario sobre las tecnologías utilizadas.

\section{Datos $LN(0,1)$}

En esta sección se analizan los datos generados utilizando que los ingresos de una población siguen una distribución $LogNormal(0,1)$. Se estudiará las distribuciones de las mediciones con distinta granularidad así como la correlación entre las mismas.

\section{Datos $LN(0,\sigma^2)$}

En esta sección se analizan los datos generados utilizando una distribución $LogNormal(0,\sigma^2)$ para distintos valores de $\sigma^2$. Trabajar con distintos valores de $\sigma^2$ permite una mayor variedad y riqueza de análisis.

\subsection{Comportamiento condicional}

Se estudia el comportamiento condicional a $\sigma^2$, tratándolo como 10 poblaciones independientes y estudiándolas.

\subsection{Interacción con granularidad}

Se estudia como se comportan las mediciones con distintas granularidades al modificar $\sigma^2$

\subsection{Comportamiento no condicional}

Se trabajará con toda la muestra sin condicionar por $\sigma^2$, es decir tratándola como una unica población heterogénea.

\chapter{Datos reales: Distribución de los ingresos}

En esta sección se utilizán datos reales de distribución del ingreso de distintos países y años para estudiar su distribución y compararlos con los datos sintéticos antes generados.

\section{Datos disponibles}

Se presenta la base de datos disponible. Se comentan los distintos alcances (nacional, urbano y rural) y tipos (ingresos y consumo) de las encuestas disponibles.

\section{Distribución empírica}

Se estudian las distribuciones empíricas disponibles y se las compara con los datos sintéticos generados en el capítulo anterior.

\chapter{Datos reales: Mejora de la capacidad predictiva}

En este capítulo exploramos si las modificaciones propuestas en la medición del bienestar económico mejoran la capacidad de predecir otros indicadores socioeconómicos.

\section{Esperanza de Vida al Nacer (EVN)}

Se estudia pormenorizadamente el caso de la EVN, discutiendo distintas formas de evaluar la capacidad predictiva de las distintas mediciones. Se utiliza el año como variable de control para evitar que una correlación entre el bienestar económico y el año y entre el año y la EVN aumente artificialmente la correlación entre el bienestar económico y la EVN.


\section{Otros indicadores}

Utilizando la métrica comparativa finalmente desarrollada para el caso de la EVN, se aplica la misma métrica a otros indicadores socioeconómicos (ej: Alfabetización a los 15 años) para ver cómo se comportan en estos casos.



\chapter{Conclusiones}

Título autoexplicativo

\chapter{Posibles trabajos futuros}

Título autoexplicativo

\chapter{Citas bibliográficas}

Título autoexplicativo



\end{document}